\subsubsection{真正 arity 维度的 primitive 素数谱：Dirichlet 分层与 Mertens 常数张量}\label{subsec:arity-dirichlet-mertens-tensor}
本小节把 40 状态真实输入核（附录 \ref{app:real-input-40-abel}）的 primitive 轨道按三个可加观测量做 Dirichlet 分层，并把每个同余类压缩成一个 Abel/Hadamard 有限部常数项（Mertens 常数漂移）。该输出是“素数影子”的下一层硬量化：既给出同余类的常数张量，也给出等分布收敛的最坏误差指数。

\paragraph{三条分层量（arity-charge + 两类隐藏/触发维度）}
对一条 primitive 轨道 $\gamma$，记其沿边的三类可加计数为
$$
\Phi_e(\gamma),\qquad \Phi_-(\gamma),\qquad \Phi_2(\gamma),
$$
它们分别对应附录 \ref{app:real-input-40-logm-theta} 中的三条边势 $\varphi_e,\varphi_-,\varphi_2$ 的轨道和。
在本文选择的三维 Dirichlet 分层中，我们以 $\chi$ 作为 arity-charge，并以 $N_-$ 与 $N_e$ 作为两条事件计数轴（其中 $N_e$ 是输出位计数；$N_2$ 仅作为构造 $\chi$ 的辅助量保留在定义中）。
定义轨道层的 arity-charge 与事件计数
$$
\chi(\gamma):=\Phi_e(\gamma)-\Phi_2(\gamma)\in\ZZ,
\qquad
N_-(\gamma):=\Phi_-(\gamma)\in\ZZ,
\qquad
N_e(\gamma):=\Phi_e(\gamma)\in\ZZ,
\qquad
N_2(\gamma):=\Phi_2(\gamma)\in\ZZ.
$$
在“真实输入 40 状态核”的同步边标记 $(s,d)\mapsto(e,t)$ 口径下，可取每步观测量
$$
\chi:=e-\mathbf{1}_{\{d=2\}}\in\{-1,0,1\},
\qquad
\nu:=\mathbf{1}_{\{t\in Q_-\}}\in\{0,1\},
\qquad
e\in\{0,1\},
\qquad
\xi:=\mathbf{1}_{\{d=2\}}\in\{0,1\}.
$$
从而
$$
\Phi_e(\gamma)=\sum_{k=1}^{|\gamma|} e_k,
\qquad
\Phi_-(\gamma)=\sum_{k=1}^{|\gamma|} \nu_k,
\qquad
\Phi_2(\gamma)=\sum_{k=1}^{|\gamma|} \xi_k,
$$
并因此 $\chi(\gamma)=\sum_{k=1}^{|\gamma|}\chi_k$、$N_-(\gamma)=\sum_{k=1}^{|\gamma|}\nu_k$、$N_e(\gamma)=\sum_{k=1}^{|\gamma|}e_k$、$N_2(\gamma)=\sum_{k=1}^{|\gamma|}\xi_k$。
其中 $\chi(\gamma)$ 是真正的 arity-charge（把“二元碰撞如何被吐出/隐藏”的净电荷编码成可加不变量），$N_-(\gamma)$ 记录进入负携带区的事件计数，而 $N_e(\gamma)$ 记录输出位计数。主增长率取真实输入核的 Perron--Frobenius 主特征值
$$
\lambda=\rho(M)=\varphi^2\approx 2.61803398875.
$$

\paragraph{多可观测量热力系统：向量势、rotation polytope 与可计算相图}
上述三条观测量不仅能做 Dirichlet 分层，也能被统一视为一个\textbf{向量值局部常势}。把 40 状态核的核心图（essential SCC；见命题 \ref{prop:real-input-40-essential-20}）记为 $\Gamma$，并对每条有向边 $e:i\to j$ 赋一个向量值观测量
$$
\Phi(e):=(\chi(e),\nu(e),\xi(e))\in\ZZ^3.
$$
对任意参数向量 $\mathbf{u}=(u_\chi,u_-,u_2)\in(\CC^\times)^3$，可定义多参权重
$$
w_{\mathbf{u}}(e):=\mathbf{u}^{\Phi(e)}:=u_\chi^{\chi(e)}u_-^{\nu(e)}u_2^{\xi(e)},
$$
并得到多参带权邻接矩阵族
$$
M(\mathbf{u})_{ij}:=\sum_{e:i\to j} w_{\mathbf{u}}(e),
$$
其 $\zeta$--行列式与 primitive Euler 乘积仍按 SFT 口径成立（\cite{BrinStuck2002,ParryPollicott1990Zeta}）。本文前述的三维扭曲矩阵 $M_{j_1,j_2,j_3}$ 正是取 $\mathbf{u}$ 为有限群角色（单位圆上的根）时的特例。

更进一步，对边移位系统（SFT）与向量势 $\Phi$，定义 rotation set（旋转集合）
$$
\mathrm{Rot}(\Phi)
:=
\left\{\int \Phi\,d\mu:\ \mu\ \text{是}\ \Gamma\ \text{上的}\ \sigma\text{-不变概率测度}\right\}\subset\RR^3.
$$
当 $\Gamma$ 为传递 SFT 且 $\Phi$ 仅依赖长度 $2$ 的 cylinder（即仅依赖边标签）时，$\mathrm{Rot}(\Phi)$ 为\textbf{凸多面体}，周期点（周期轨道）的 rotation vectors 在其中稠密，且每个内点由遍历测度实现（\cite{Ziemian1995RotationSFT}）。因此，“隐含属性/操作类型配比”的长期平均并不是直觉对象，而是一个\textbf{有限维凸多面体坐标}。

rotation polytope 的支撑函数可通过\textbf{最小平均回路权}精确计算：对任意方向 $\theta\in\RR^3$，设标量化势
$$
\phi_\theta(e):=\langle\theta,\Phi(e)\rangle,
$$
则有（见 \cite{Karp1978MinCycleMean} 的最小平均回路刻画）
$$
h_{\min}(\theta):=\min_{v\in\mathrm{Rot}(\Phi)}\langle\theta,v\rangle
=\min_{\gamma}\frac{1}{|\gamma|}\sum_{e\in\gamma}\phi_\theta(e),
$$
右端对所有周期轨道（回路） $\gamma$ 取最小。由此可把“纯相”定义为被某个方向 $\theta$ 暴露的面；$\theta$-空间中使同一面被选中的区域构成 rotation polytope 的 normal fan，从而得到一个\textbf{有限个相的核相图}。这一视角与局部常势的零温极限（$\beta\to+\infty$）一致：零温选择集中在有限并的周期/子移位分量上（参见 \cite{ChazottesGambaudoUgalde2011ZeroTempLocConst} 的一般理论）。

作为可审计示例，我们在 essential 20 状态核心上取 $\Phi=(\chi,\nu,\xi)$，对若干方向 $\theta$ 做标量化并在图上求最小平均回路权与见证周期轨道，输出其周期均值（rotation vector）：
\input{sections/generated/tab_real_input_40_rotation_polytope_sample}

\begin{remark}[无回溯素数影子与覆盖的 \texorpdfstring{$L$}{L}-函数分解]\label{rem:rotation-polytope-bartholdi}
上文的 primitive Euler 乘积基于“周期点/闭路”口径；但在“原子更新近可逆”的核语境里，更自然的素数影子常需\textbf{排除回溯}（做了又撤销）。这对应 Ihara--Bass--Hashimoto/Bartholdi 的无回溯素回路框架：对（有向/带权）图可定义加权 Bartholdi $\zeta$ 与相应的加权 $L$-函数，并给出行列式表达；在图覆盖（有限群标记）下还存在 Artin 型分解：覆盖的 $\zeta$ 可分解为各不可约表示的扭曲 $L$-函数乘积（见 \cite{ChoeKwakParkSato2007Bartholdi,Sato2008WeightedBartholdiCoverings}）。与本节的 Dirichlet 扭曲谱半径比 $\rho/\lambda$ 类似，覆盖/扭曲的谱半径比提供了“同余类偏差”的可计算指数尺度；在 Axiom A/超曲性动力系统中，这一类比进一步对应到 Chebotarev 型分布定理（参见 \cite{PollicottSharp2007Chebotarev}）。
\end{remark}

\paragraph{\texorpdfstring{$((2,2,2))$}{(2,2,2)}：arity--Mertens 常数张量}
记 residue 向量
$$
(r_1,r_2,r_3)=\bigl(\chi(\gamma)\bmod 2,\ N_-(\gamma)\bmod 2,\ N_e(\gamma)\bmod 2\bigr)\in(\ZZ/2)^3.
$$
令 $p_{n,r_1r_2r_3}$ 表示长度为 $n$ 的 primitive 轨道中落在该 residue 的轨道数（按轨道计）。定义 Mertens 部分和
$$
S_{r_1r_2r_3}(N):=\sum_{n\le N}\frac{p_{n,r_1r_2r_3}}{\lambda^n}.
$$
则存在常数 $C_{r_1r_2r_3}$ 使得
$$
S_{r_1r_2r_3}(N)=\frac{1}{8}\log N + C_{r_1r_2r_3}+o(1),
\qquad
\sum_{r_1,r_2,r_3\in\{0,1\}}C_{r_1r_2r_3}=\mathsf{M},
$$
其中总 Mertens 常数满足
$$
\mathsf{M}\approx 0.322903002999.
$$
对应的常数张量数值与最坏误差指数（$\rho/\lambda$）由脚本
\texttt{scripts/exp\_sync\_kernel\_real\_input\_40\_arity\_3d.py}
自动生成并输出为可 \verb|\input| 的 \LaTeX\ 片段：
\input{sections/generated/tab_real_input_40_arity_dirichlet_mertens_222}

\paragraph{\texorpdfstring{$((3,2,2))$}{(3,2,2)}：把 arity 维度提升为真正 Dirichlet 分层}
把第一维升级为
$$
(a,b,c)=\bigl(\chi(\gamma)\bmod 3,\ N_-(\gamma)\bmod 2,\ N_e(\gamma)\bmod 2\bigr)\in(\ZZ/3)\times(\ZZ/2)\times(\ZZ/2),
$$
并令 $p_{n,a,b,c}$ 为相应同余类的 primitive 轨道计数，定义
$$
S_{a,b,c}(N):=\sum_{n\le N}\frac{p_{n,a,b,c}}{\lambda^n}.
$$
则同样有
$$
S_{a,b,c}(N)=\frac{1}{12}\log N + C_{a,b,c}+o(1),
\qquad
\sum_{a\in\{0,1,2\}}\sum_{b,c\in\{0,1\}}C_{a,b,c}=\mathsf{M}.
$$
对应的常数张量数值与最坏误差指数（$\rho/\lambda$）同样由脚本
\texttt{scripts/exp\_sync\_kernel\_real\_input\_40\_arity\_3d.py}
自动生成并输出为可 \verb|\input| 的 \LaTeX\ 片段：
\input{sections/generated/tab_real_input_40_arity_dirichlet_mertens_322}

\paragraph{\texorpdfstring{$((3,3,3))$}{(3,3,3)}：三维 Dirichlet--Mertens 常数张量}
进一步把三条观测量同时提升到模 $3$ 的 Dirichlet 分层：
$$
(a,b,c)=\bigl(\chi(\gamma)\bmod 3,\ N_-(\gamma)\bmod 3,\ N_e(\gamma)\bmod 3\bigr)\in(\ZZ/3)^3.
$$
令 $p_{n,a,b,c}$ 表示长度为 $n$ 的 primitive 轨道中落在该 residue 的轨道数（按轨道计），并定义
$$
S_{a,b,c}(N):=\sum_{n\le N}\frac{p_{n,a,b,c}}{\lambda^n}.
$$
则存在常数张量 $C^{(3,3,3)}_{a,b,c}$ 使得
$$
S_{a,b,c}(N)=\frac{1}{27}\log N + C^{(3,3,3)}_{a,b,c}+o(1),
\qquad
\sum_{a,b,c\in\{0,1,2\}} C^{(3,3,3)}_{a,b,c}=\mathsf{M}.
$$
对应的常数张量数值（含一维边际常数、以及最坏扭曲角色索引）由脚本
\texttt{scripts/exp\_sync\_kernel\_real\_input\_40\_arity\_3d.py}
直接生成并输出为可 \verb|\input| 的 \LaTeX\ 片段：
% AUTO-GENERATED by scripts/exp_sync_kernel_real_input_40_arity_3d.py
% (3,3,3) Dirichlet--Mertens constants tensor for (chi mod 3, N_- mod 3, N_2 mod 3).
\paragraph{$((3,3,3))$：三维 Dirichlet--Mertens 常数张量（可复现输出）}
按 $a=\chi(\gamma)\bmod 3$ 分三张 $3\times 3$ 切片（行 $b=N_-(\gamma)\bmod 3$，列 $c=N_2(\gamma)\bmod 3$），常数项为：
$$
a=0:\quad
\begin{pmatrix}
\phantom{-}0.51457124 & -0.10167814 & -0.09054112\\
-0.11494313 & -0.04637293 & -0.07761890\\
-0.11505981 & \phantom{-}0.10734114 & -0.08930265
\end{pmatrix}
$$
$$
a=1:\quad
\begin{pmatrix}
\phantom{-}0.92442842 & -0.05868184 & -0.08110602\\
\phantom{-}0.15485275 & \phantom{-}0.00403884 & -0.08650366\\
-0.11349611 & -0.02935305 & -0.08990830
\end{pmatrix}
$$
$$
a=2:\quad
\begin{pmatrix}
\phantom{-}0.31218143 & -0.10160005 & -0.08465747\\
-0.11604115 & -0.01446314 & -0.08537890\\
-0.07159741 & -0.03006004 & -0.09614697
\end{pmatrix}
$$
\paragraph{指数级误差：最坏扭曲谱半径比（可复现输出）}
数值上
$$
\rho_{3,3,3}\approx 2.452276675453,\qquad \frac{\rho_{3,3,3}}{\lambda}\approx 0.936686340205.
$$
达到最坏谱半径的非平凡角色索引（$j=(j_1,j_2,j_3)$）为：
$$
(0,1,1),\ (0,2,2)
$$
\paragraph{一维边际：只看 $\chi\bmod 3$ 的 Dirichlet--Mertens 常数}
对 $b,c$ 求和得到 $C_a^{(3)}:=\sum_{b,c} C^{(3,3,3)}_{a,b,c}$：
$$
\bigl(C_0^{(3)},C_1^{(3)},C_2^{(3)}\bigr)\approx\bigl(-0.01360430,+0.62427101,-0.28776371\bigr).
$$


\paragraph{指数级误差：最坏扭曲谱半径比}
令 $\omega_3:=e^{2\pi i/3}$。对每个非平凡角色 $(j_1,j_2,j_3)\in(\ZZ/3)^3\setminus\{(0,0,0)\}$，定义三变量扭曲矩阵为
$$
M(\omega_3^{j_1},\omega_3^{j_2},\omega_3^{j_3})
$$
（即沿边对三条可加观测量分别施加角色权重），并令
$$
\rho_{3,3,3}
:=
\max_{(j_1,j_2,j_3)\neq(0,0,0)}
\rho\bigl(M(\omega_3^{j_1},\omega_3^{j_2},\omega_3^{j_3})\bigr).
$$
其数值输出亦包含在上述 \verb|\input| 片段中。
因此同余类趋于等分布的误差为指数级：
$$
p_{n,a,b,c}=\frac{p_n}{27}+O\!\left(\frac{\rho_{3,3,3}^n}{n}\right),
\qquad
p_n\sim\frac{\lambda^n}{n}.
$$

\paragraph{\texorpdfstring{$((3,3,5))$}{(3,3,5)}：第三维固定为 \texorpdfstring{$N_2\bmod 5$}{N2 mod 5}（arity--Dirichlet--Mertens 张量）}
由于 $\chi(\gamma)=N_e(\gamma)-N_2(\gamma)$，仅在模 $3$ 的角色扭曲层面，$\chi(\gamma)\bmod 3$ 与 $N_e(\gamma)\bmod 3$ 之间会出现退化的“置换等价”。为得到一个对碰撞触发机制（$d=2$）真正敏感、且在工程上可审计的三维常数张量，我们固定第三维为二元碰撞触发计数
\(
N_2(\gamma)\bmod 5
\)：
$$
(a,b,c)=\bigl(\chi(\gamma)\bmod 3,\ N_-(\gamma)\bmod 3,\ N_2(\gamma)\bmod 5\bigr)\in(\ZZ/3)\times(\ZZ/3)\times(\ZZ/5).
$$
相应地，若 $p_{n,a,b,c}$ 为同余类的 primitive 轨道数（按轨道计），并令
$$
S_{a,b,c}(N):=\sum_{n\le N}\frac{p_{n,a,b,c}}{\lambda^n},
$$
则存在常数张量 $C^{(3,3,5)}_{a,b,c}$ 使得
$$
S_{a,b,c}(N)=\frac{1}{45}\log N + C^{(3,3,5)}_{a,b,c}+o(1),
\qquad
\sum_{a=0}^2\sum_{b=0}^2\sum_{c=0}^4 C^{(3,3,5)}_{a,b,c}=\mathsf{M}.
$$
\begin{proposition}[为何 \texorpdfstring{$N_2\bmod 5$}{N2 mod 5} 必然产生新结构（非置换等价）]\label{prop:arity-335-n2-nondegenerate}
在本节的三维分层
$$
(a,b,c)=\bigl(\chi(\gamma)\bmod 3,\ N_-(\gamma)\bmod 3,\ N_2(\gamma)\bmod 5\bigr),
$$
第三维 $c=N_2(\gamma)\bmod 5$ 不能被前两维（特别是 $\chi(\gamma)\bmod 3$）所消去；因此对固定的 $(a,b)\in(\ZZ/3)^2$，primitive 轨道会沿 $c\in\ZZ/5$ 发生不可约分裂（不是重标记置换）。
\end{proposition}

\begin{proof}
该核上恒有一步关系 $\chi=e-\xi$，从而轨道层恒等式
$$
\chi(\gamma)=N_e(\gamma)-N_2(\gamma),
$$
但本分层只保留 $\chi(\gamma)\bmod 3$ 而不保留 $N_e(\gamma)\bmod 5$。因此即使固定 $a=\chi(\gamma)\bmod 3$，也无法由 $a$ 决定 $N_2(\gamma)\bmod 5$（等价地，无法由 $\chi(\gamma)\bmod 3$ 决定 $\chi(\gamma)\bmod 5$）。
换言之，第三维的 $\ZZ/5$ 信息不可能被前两维的 $\ZZ/3$ 信息“线性还原”，从而 $c$ 维度引入了真正的新自由度。
\end{proof}

\begin{proposition}[collision-cocycle 主导的慢模阻塞（\texorpdfstring{$((3,3,5))$}{(3,3,5)}；第三轴 \texorpdfstring{$N_2\bmod 5$}{N2 mod 5}）]\label{prop:arity-335-collision-dominant}
在 \(((3,3,5))\) 的三维扭曲设定下（第三轴取 $c=N_2(\gamma)\bmod 5$，扭曲角色索引 $j=(j_1,j_2,j_3)\in(\ZZ/3)^2\times(\ZZ/5)$，并令
\(
\rho(j):=\rho\bigl(M_{j_1,j_2,j_3}\bigr)
\)
为对应扭曲矩阵的谱半径），记最坏扭曲谱半径为
$$
\rho_{3,3,5}:=\max_{j\neq (0,0,0)}\rho(j).
$$
则该最大值由\textbf{纯碰撞扭曲}实现：存在 $j_3\in\{1,4\}$ 使得
$$
\rho_{3,3,5}=\rho(0,0,j_3),
$$
并且慢模阻塞完全来自 $\ZZ/5$-值的碰撞 cocycle
\(
\xi=\mathbf{1}_{\{d=2\}}
\)
（而非来自 $\chi(\gamma)\bmod 3$ 或 $N_-(\gamma)\bmod 3$ 的扭曲）。
\end{proposition}

\begin{proof}
按定义，$\rho_{3,3,5}$ 是对全部非平凡角色索引的谱半径最大值。由本节给出的可复现输出片段
\verb|\input{sections/generated/tab_real_input_40_arity_dirichlet_mertens_335_N2}|
中“达到最坏谱半径的非平凡角色索引”为 $(0,0,1),(0,0,4)$，故最大值确由 $(j_1,j_2)=(0,0)$ 的纯第三轴扭曲达到，即
\(
\rho_{3,3,5}=\rho(0,0,1)=\rho(0,0,4)
\)。
由于 $j_1=j_2=0$，该扭曲对 $\chi(\gamma)\bmod 3$ 与 $N_-(\gamma)\bmod 3$ 不施加任何角色权重，唯一的扭曲来自第三轴（即 $\xi=\mathbf{1}_{\{d=2\}}$ 的 $\ZZ/5$ 角色），从而慢模阻塞由碰撞 cocycle 单独主导。
\end{proof}

\begin{corollary}[等分布的主相关长度尺度由纯碰撞扭曲给出]\label{cor:arity-335-collision-mixing-scale}
设 \(\rho_{3,3,5}\) 与 \(\lambda=\rho(M)\) 如上。则 \(((3,3,5))\) 同余类计数的最坏扭曲误差以指数率
\(
(\rho_{3,3,5}/\lambda)^n
\)
衰减；因此可把“主相关长度尺度”定义为
$$
\tau_{\mathrm{mix}}:=\frac{1}{-\log(\rho_{3,3,5}/\lambda)},
\qquad
t_{1/2}:=\frac{\log 2}{-\log(\rho_{3,3,5}/\lambda)}.
$$
并且按命题 \ref{prop:arity-335-collision-dominant}，该最坏衰减模式由纯碰撞扭曲 $(0,0,1)$ 或 $(0,0,4)$ 产生。
\end{corollary}

\begin{remark}[慢模与折叠碰撞核的同源性信号]\label{rem:arity-335-rho-close-to-a2}
折叠门的二阶碰撞统计给出三态碰撞核 \(A_2\)（命题 \ref{prop:pom-s2-recurrence}），其 Perron 根为
$$
r_2=\rho(A_2)\approx 2.481194304092
$$
（表 \ref{tab:collision_kernel_A2_finite_part}）。
另一方面，对真实输入 40 状态核的 \(((3,3,5))\) 分层（第三轴为 \(N_2(\gamma)\bmod 5\)）可复现输出给出
$$
\rho_{3,3,5}\approx 2.488758321093
$$
（见下述表格片段）。
两者在数值上极其接近：
$$
\rho_{3,3,5}-r_2\approx 7.564\times 10^{-3},
\qquad
\frac{\rho_{3,3,5}}{r_2}\approx 1.00305.
$$
这提示：\(((3,3,5))\) 的最慢模态并非由输出端或负携带端的细节主导，而更像由“二阶碰撞等价关系的纤维积结构”（折叠碰撞核）在轨道层的投影所控制；换言之，\(N_2(\gamma)\bmod 5\) 的慢变量可能是折叠碰撞机制在有限模轴束上的可审计读出。
\end{remark}

\begin{remark}[Pisano--20 共振线索]\label{rem:arity-335-pisano}
二级摘要输出给出
\(
\tau_{\mathrm{mix}}\approx 19.7473
\)，数值上贴近 Fibonacci 模 $5$ 的 Pisano 周期 $\pi(5)=20$。该贴合提示 $\ZZ/5$ 轴上的慢模可能近似实现一个“Pisano--20”准旋转结构；此解释仅为动机性叙述，后续可通过更高模数的扫描检验其稳定性。
\end{remark}

\begin{remark}[分歧素数的动机性解释]\label{rem:arity-335-ramification}
解释性地看，黄金域 $\QQ(\sqrt5)$ 的判别式为 $5$，素数 $5$ 在该域中分歧（参见 \cite{Apostol1976} 的标准事实）。这为 “模 $5$ 轴呈现异常慢模/共振” 提供了一个数论动机，但不作为后续推理前提。
\end{remark}

\paragraph{结果（五张 \texorpdfstring{$3\times 3$}{3x3} 切片 \texorpdfstring{$+$}{+} 最坏误差指数 \texorpdfstring{$+$}{+} 两条审计）}
下述 \LaTeX\ 片段直接给出 $c=0,1,2,3,4$ 的 5 张切片、最坏谱半径比 $\rho_{3,3,5}/\lambda$ 与达到最坏的角色索引，并输出两条审计误差（总和校验与边际一致性）：
\paragraph{计算闭环（可审计）：三维扭曲 $\Rightarrow$ 扭曲常数 $\Rightarrow$ 傅里叶反演}
令 $\omega_3=e^{2\pi i/3},\ \omega_5=e^{2\pi i/5}$。对每个角色索引 $(j_1,j_2,j_3)\in\{0,1,2\}^2\times\{0,1,2,3,4\}$，在真实输入 40 状态核上构造扭曲矩阵
$$
M_{j_1,j_2,j_3}:\quad
e\mapsto (\omega_3^{j_1})^{\chi(e)}(\omega_3^{j_2})^{\nu(e)}(\omega_5^{j_3})^{\xi(e)},
$$
其中 $\xi=\mathbf{1}_{\{d=2\}}$、$\chi=e-\xi$、$\nu=\mathbf{1}_{\{s'\in Q_-\}}$ 均为一步可加量（定义见附录 \ref{app:real-input-40-logm-theta} 与 \ref{app:real-input-40-dirichlet}）。
按 SFT 的 $\zeta$--行列式口径（\cite{BrinStuck2002}）定义
$$
\zeta_{j_1,j_2,j_3}(z)=\frac{1}{\det(I-zM_{j_1,j_2,j_3})},
\qquad
\log\zeta_{j_1,j_2,j_3}(z)=\sum_{n\ge 1}\frac{\Tr(M_{j_1,j_2,j_3}^n)}{n}z^n.
$$
对非平凡角色 $(j_1,j_2,j_3)\neq(0,0,0)$，在 $z=\lambda^{-1}$ 处用 Möbius--$\log\zeta$ 的收敛恒等式得到扭曲常数（见附录 \ref{app:real-input-40-finite-part} 的同型推导）：
$$
C(j_1,j_2,j_3):=\sum_{m\ge 1}\frac{\mu(m)}{m}\log\zeta_{j_1m,\ j_2m,\ j_3m}(\lambda^{-m}),
$$
其中 $(j_1m,\ j_2m,\ j_3m)$ 按模 $(3,3,5)$ 取值。最后由三维离散傅里叶反演得到常数张量：
$$
C^{(3,3,5)}_{a,b,c}
=\frac{\mathsf{M}}{45}
\frac{1}{45}\!\!\!\sum_{\substack{0\le j_1\le 2\\0\le j_2\le 2\\0\le j_3\le 4\\(j_1,j_2,j_3)\ne(0,0,0)}}
\omega_3^{-j_1 a}\,\omega_3^{-j_2 b}\,\omega_5^{-j_3 c}\ C(j_1,j_2,j_3).
$$
该闭环同时输出最坏扭曲谱半径比 $\rho_{3,3,5}/\lambda$（指数误差尺度）与两条校验：
$$
\sum_{a,b,c}C^{(3,3,5)}_{a,b,c}=\mathsf{M},
\qquad
\sum_{c=0}^{4}C^{(3,3,5)}_{a,b,c}=C^{(3,3)}_{a,b}.
$$

常数张量（按 $c=N_2(\gamma)\bmod 5$ 给出 5 张 $3\times 3$ 切片）、最坏扭曲谱半径比 $\rho_{3,3,5}/\lambda$、以及上述校验的数值误差，由脚本
\texttt{scripts/exp\_sync\_kernel\_real\_input\_40\_arity\_3d.py}
直接生成（\texttt{--third-axis N2}）并输出为可 \verb|\input| 的 \LaTeX\ 片段：
\input{sections/generated/tab_real_input_40_arity_dirichlet_mertens_335_N2}

\paragraph{二级推论（由上述五切片与最坏谱半径比直接推出）}
为了把“表格输出”进一步压缩为可直接引用的结构性结论，我们额外输出一份仅由上述五切片与最坏谱半径比计算得到的二级摘要（同一脚本生成）：
% AUTO-GENERATED by scripts/exp_sync_kernel_real_input_40_arity_3d.py
% Derived consequences for the (3,3,5) tensor with third axis N2 mod 5.
\paragraph{从 $((3,3,5))$ 输出推导出的二级结论（可复现）}
下述量全部由上表的 5 张 $3\times 3$ 切片与最坏扭曲谱半径比直接计算得到。
\paragraph{（i）慢模完全由碰撞 cocycle 主导}
最坏谱半径比与达到最坏的角色索引为：
$$
\frac{\rho_{3,3,5}}{\lambda}\approx 0.950621088873.
$$
$$
(0,0,1),\ (0,0,4)
$$
由此定义主相关长度尺度（指数衰减的时间常数）与半衰期：
$$
\tau_{\mathrm{mix}}:=\frac{1}{-\log(\rho_{3,3,5}/\lambda)}\approx 19.747340579191,\qquad t_{1/2}:=\frac{\log 2}{-\log(\rho_{3,3,5}/\lambda)}\approx 13.687813446023.
$$
\paragraph{（ii）碰撞同余类 $c$ 的常数漂移（总偏置）}
定义 $S_c:=\sum_{a=0}^2\sum_{b=0}^2 C^{(3,3,5)}_{a,b,c}$。数值上
$$
(S_0,S_1,S_2,S_3,S_4)\approx(+1.685330,+0.130189,-0.321067,-0.528302,-0.643249).
$$
\paragraph{（iii）三维耦合在 $c$ 轴上的局部化（切片方差）}
定义每张切片在 $3\times 3$ 网格上的标准差（按 $9$ 个条目平均）：
$$
\bigl(\mathrm{sd}(c=0),\dots,\mathrm{sd}(c=4)\bigr)\approx\bigl(0.347585,0.0591959,0.00515942,0.000787443,0.000137309\bigr).
$$
\paragraph{（iv）$\mathbb{{Z}}/5$ 轴的傅里叶压缩指纹（均值 + 两个复模）}
对每个固定的 $(a,b)$，令 $f_{a,b}(c):=C^{(3,3,5)}_{a,b,c}$（$c\in\ZZ/5$）。其离散傅里叶系数为
$$
\widehat f_{a,b}(j):=\frac15\sum_{c=0}^{4} f_{a,b}(c)\,\omega_5^{-jc},\qquad \omega_5=e^{2\pi i/5}.
$$
从而有恒等分解（对所有 $c$ 精确成立）
$$
f_{a,b}(c)=\widehat f_{a,b}(0)+2\Re\bigl(\widehat f_{a,b}(1)\,\omega_5^{c}\bigr)+2\Re\bigl(\widehat f_{a,b}(2)\,\omega_5^{2c}\bigr).
$$
记 $\overline C_{b,a}:=\widehat f_{a,b}(0)$（行 $b$、列 $a$）与两张复矩阵 $A^{(1)}_{b,a}:=\widehat f_{a,b}(1)$、$A^{(2)}_{b,a}:=\widehat f_{a,b}(2)$。则：
$$
\overline C=
\begin{pmatrix}
\phantom{-}0.06447040 & \phantom{-}0.15692811 & \phantom{-}0.02518478\\
-0.04778699 & \phantom{-}0.01447759 & -0.04317664\\
-0.01940426 & -0.04655149 & -0.03956088
\end{pmatrix}
$$
并且（以条目绝对值的最大值度量）
$$
\max_{a,b}|A^{(1)}_{a,b}|\approx 0.201317996544,\qquad\max_{a,b}|A^{(2)}_{a,b}|\approx 0.200144063931,\qquad\frac{\max|A^{(2)}|}{\max|A^{(1)}|}\approx 0.994168764674.
$$
其中 $A^{(1)}$ 与 $A^{(2)}$ 的具体条目可由同一脚本在 JSON 输出中复核。


\begin{corollary}[碰撞同余类的常数级偏置（\texorpdfstring{$S_c$}{Sc} 分解）]\label{cor:arity-335-collision-residue-drift}
令
\[
S_c:=\sum_{a=0}^{2}\sum_{b=0}^{2}C^{(3,3,5)}_{a,b,c},
\qquad c\in\ZZ/5.
\]
则在真实输入 40 状态核的 \(((3,3,5))\) 分层（第三轴取 $c=N_2(\gamma)\bmod 5$）下，可复现输出显示：$S_0$ 显著为正，而 $S_2,S_3,S_4$ 为负，并且 $\sum_{c=0}^4 S_c=\mathsf{M}$。因此在 Abel/Hadamard 有限部（常数级漂移）口径下，primitive 轨道对碰撞同余类 $c$ 的偏置主要体现为“$c=0$ 的强正漂移 + 其余同余类的系统性负漂移”。
\end{corollary}

\begin{corollary}[符号相变阈值与正漂移稀疏性]\label{cor:arity-335-sign-sparsity}
对 \(((3,3,5))\) 的五张切片逐项检查可复现输出显示：
\[
C^{(3,3,5)}_{a,b,c}>0\ \Longrightarrow\ c\in\{0,1\},
\]
并且当 $c=2,3,4$ 时全部 $9$ 个格点严格为负。由此正漂移仅出现在 $c=0,1$ 两片中。进一步统计得到
\[
\#\{C>0\}=9,\qquad \#\{C<0\}=36,
\]
以及
\[
\sum_{C>0}C\approx 2.303140,\qquad
\sum_{C<0}C\approx -1.980237,\qquad
\sum_{a,b,c}C^{(3,3,5)}_{a,b,c}=\mathsf{M}.
\]
因此总常数 $\mathsf{M}$ 是“大正漂移”与“大负漂移”的强抵消残差，而非单侧累积。
\end{corollary}

\begin{remark}[抵消因子与共振胞元]\label{rem:arity-335-cancel-resonant}
定义抵消因子
\[
\kappa_{\mathrm{cancel}}:=\frac{\sum_{C>0}C}{\sum_{a,b,c}C^{(3,3,5)}_{a,b,c}}
=\frac{\sum_{C>0}C}{\mathsf{M}},
\]
则在 \(((3,3,5))\) 输出上
\(
\kappa_{\mathrm{cancel}}\approx 7.13
\)，显示 $\mathsf{M}$ 是强正负漂移抵消后的残差。
此外，由表格可复现输出可见最大单元发生在
\(
(a,b,c)=(1,0,0)
\)，其值
\(
C^{(3,3,5)}_{1,0,0}\approx 0.95882121
\)
已约为总常数 $\mathsf{M}$ 的 $2.97$ 倍，体现“共振胞元”对常数项的主导贡献。
\end{remark}

\begin{corollary}[三维耦合对 \texorpdfstring{$(a,b)$}{(a,b)} 的依赖在 $c$ 轴上强局部化]\label{cor:arity-335-localized-ab-coupling}
把每张 $c$-切片看作 $3\times 3$ 矩阵（行 $b$，列 $a$），并令 $\mathrm{sd}(c)$ 为该矩阵在 $9$ 个条目上的标准差（按条目平均）。则可复现输出显示：$\mathrm{sd}(c)$ 随 $c$ 增大呈数量级下降，且在 $c\ge 2$ 时切片已接近常数矩阵（对 $(a,b)$ 的依赖显著减弱），在 $c=3,4$ 时更是几乎完全消失。故在常数级漂移分辨率下，\(((3,3,5))\) 的非平凡三维耦合主要集中在少数碰撞同余类（特别是 $c=0$ 与 $c=1$）上，而尾部同余类呈近因子化行为。
\end{corollary}

\begin{corollary}[三相分层：$c=0$ 的结构尖峰、$c=1$ 的过渡与 $c\ge 2$ 的近因子化尾部]\label{cor:arity-335-phase-stratification}
以二级摘要输出中的 $\max\mathrm{dev}(c)$ 与 $\Delta_c$ 为诊断量，则 \(((3,3,5))\) 常数张量沿碰撞同余类轴 $c\in\ZZ/5$ 呈现明显的三相分层：
\begin{itemize}
  \item \textbf{尖峰相（$c=0$）}：切片内在 $(a,b)$ 上高度非均匀；
  \item \textbf{过渡相（$c=1$）}：切片非均匀度显著下降但仍可见结构；
  \item \textbf{尾部相（$c\ge 2$）}：切片接近 rank-$0$ 常数矩阵（近因子化）。
\end{itemize}
这比命题 \ref{prop:arity-335-n2-nondegenerate} 的“非置换等价”更强：不仅第三轴引入了新自由度，而且该自由度的结构性主要局部化在小同余类上。
\end{corollary}

\begin{remark}[尾部近标量化指纹 \texorpdfstring{$\eta_{\mathrm{tail}}$}{eta_tail}]\label{rem:arity-335-eta-tail}
把 $c$-切片记为矩阵 \(C_c:=(C^{(3,3,5)}_{a,b,c})_{a,b}\)，并记其均值 \(\overline C_c:=\frac19\sum_{a,b}C^{(3,3,5)}_{a,b,c}\)。
定义尾部近标量化指纹
$$
\eta_{\mathrm{tail}}
:=\max_{c\in\{2,3,4\}}
\frac{\|C_c-\overline C_c\,\mathbf{1}\mathbf{1}^{\top}\|_F}{|\overline C_c|}.
$$
该量是无量纲且可审计的：它直接度量“尾部切片偏离 rank-$0$ 常数矩阵”的相对强度，并比仅报告 $\mathrm{sd}(c)$ 更稳定（已做尺度归一）。由二级摘要输出可复现得到 \(\eta_{\mathrm{tail}}\approx 0.433880\)。
\end{remark}

\begin{remark}[尾部近标量化与有理尾律迹象]\label{rem:arity-335-rational-tail}
令 $\overline C_c:=\frac19\sum_{a,b}C^{(3,3,5)}_{a,b,c}$。由二级摘要输出可得
\[
\overline C_2\approx -0.035674,\qquad
\overline C_3\approx -0.058700,\qquad
\overline C_4\approx -0.071472,
\]
并且 $\max\mathrm{dev}(c)$ 在 $c=3,4$ 已降到 $10^{-3}$ 量级，故尾部切片几乎为标量矩阵。数值上它们分别接近
\(
-(1/28),\ -(1/17),\ -(1/14)
\)
的简单倒数（相对误差约 $10^{-3}$），提示存在可检验的“有理尾律”数值刚性。
\end{remark}

\begin{corollary}[碰撞稀少但主导慢模与常数各向异性（near-resonant cocycle 诊断）]\label{cor:arity-335-rare-collision-dominant}
在真实输入 40 状态核的三势压力 $P(\theta_e,\theta_-,\theta_2)=\log\rho(M_\theta)$ 下，碰撞触发 $\varphi_2=\mathbf{1}_{\{d=2\}}$ 的典型密度为
\[
\partial_{\theta_2}P(0)=\frac{3-\sqrt5}{10}\approx 0.076393,
\]
即在最大熵基准附近碰撞事件是稀少的（见上文关于 $\nabla P(0)$ 的闭式）。
但对 \(((3,3,5))\) 的 Dirichlet--Mertens 分层（第三轴为 $c=N_2(\gamma)\bmod 5$），可复现输出同时显示：
\begin{itemize}
  \item \textbf{慢模主导：}最坏扭曲谱半径由纯碰撞扭曲达到（命题 \ref{prop:arity-335-collision-dominant}），其衰减尺度由推论 \ref{cor:arity-335-collision-mixing-scale} 给出；
  \item \textbf{常数各向异性：}聚合漂移向量 $S_c=\sum_{a,b}C^{(3,3,5)}_{a,b,c}$ 在 $c$ 轴上呈强非均匀分配（推论 \ref{cor:arity-335-collision-residue-drift} 与二级摘要输出）。
\end{itemize}
因此尽管碰撞事件的典型密度很低，它仍能同时决定最坏等分布误差的指数尺度与常数级漂移的主要方向；这正是“近共边界/近旋转因子”的经验信号：稀少事件携带相干相位，使得其 \(\ZZ/5\)-cocycle 在统计与素数影子两端同时成为瓶颈。
\end{corollary}

\begin{proposition}[\texorpdfstring{$\ZZ/5$}{Z/5} 轴的傅里叶压缩指纹（精确恒等分解）]\label{prop:arity-335-fourier-fingerprint}
对每个固定的 $(a,b)$，令 $f_{a,b}(c):=C^{(3,3,5)}_{a,b,c}$（$c\in\ZZ/5$），并记 $\omega_5=e^{2\pi i/5}$。定义离散傅里叶系数
\[
\widehat f_{a,b}(j):=\frac15\sum_{c=0}^{4} f_{a,b}(c)\,\omega_5^{-jc},
\qquad j\in\{0,1,2,3,4\}.
\]
则对所有 $c$ 有恒等分解
\[
f_{a,b}(c)
=\widehat f_{a,b}(0)
+2\Re\bigl(\widehat f_{a,b}(1)\,\omega_5^{c}\bigr)
+2\Re\bigl(\widehat f_{a,b}(2)\,\omega_5^{2c}\bigr),
\]
从而五张切片在 $c$ 轴上的全部信息可被等价地压缩为三张矩阵：均值矩阵 $\widehat f(0)$ 与两张复矩阵 $\widehat f(1),\widehat f(2)$。该分解与其数值（含 $\|\widehat f(1)\|_\infty,\|\widehat f(2)\|_\infty$）由同一脚本输出并可审计复核。
\end{proposition}

\begin{proposition}[低秩可压缩：\texorpdfstring{$((3,3,5))$}{(3,3,5)} 张量的三模近似指纹]\label{prop:arity-335-low-rank}
把常数张量 $C^{(3,3,5)}_{a,b,c}$ 视为实矩阵 $T\in\RR^{5\times 9}$（行指标为 $c\in\{0,1,2,3,4\}$，列指标为扁平化后的 $(a,b)\in(\ZZ/3)^2$）。令其奇异值为 $s_1\ge \cdots\ge s_5$。二级摘要输出显示
\[
\frac{s_1^2}{\sum s_i^2}\approx 0.935619,\qquad
\frac{s_1^2+s_2^2}{\sum s_i^2}\approx 0.986715,\qquad
\frac{s_1^2+s_2^2+s_3^2}{\sum s_i^2}\approx 0.999866,
\]
因此仅用前三个奇异模态即可捕获几乎全部 Frobenius 能量，整个 $45$ 项张量在常数级分辨率下可被极低维近似稳定表征（可作为新的“核指纹”不变量）。
\end{proposition}

\begin{remark}[near-coboundary 二次律：对更高模数的可检验预测]\label{rem:arity-335-near-coboundary-scaling}
由二级摘要输出的缺口 $1-\rho_{3,3,5}/\lambda$ 可估计一个“近共边界强度”常数 $\kappa$，并给出
\[
\rho_{3,3,p}/\lambda\approx 1-\kappa(2\pi/p)^2
\]
的可检验预测（例如 $p\in\{7,11,13\}$）。该二次律是“简单 PF 本征值的解析微扰”在小角扭曲下的标准形状；若后续计算显示 $p=5$ 明显偏离该曲线，则可把它视为真正的特异共振模数。
\end{remark}

\begin{proposition}[小角扭曲的二次律：$\kappa$ 的方差密度识别（可证明/可算）]\label{prop:arity-335-kappa-variance}
取碰撞 cocycle \(\xi=\mathbf{1}_{\{d=2\}}\)（沿边取值 \(0/1\)），并令三轴扭曲中第三轴取 \(N_2(\gamma)\bmod p\)。
对 \(\theta\in\RR\) 定义纯碰撞扭曲矩阵
$$
M_\theta:=M(1,1,e^{\mathrm{i}\theta})\qquad(\text{仅对 }\xi\text{ 施加相位}).
$$
记 \(\Lambda(\theta)\) 为 \(M_\theta\) 的主特征值分支（\(\Lambda(0)=\lambda\)）。
由于本系统为有限维原始非负矩阵的解析扰动，\(\Lambda(\theta)\) 在 \(\theta=0\) 的邻域解析（参见一般的有限维谱扰动理论）。
则存在常数 \(\mu_\xi\in\RR\) 与 \(\sigma_\xi^2\ge 0\) 使得
$$
\log\frac{\Lambda(\theta)}{\lambda}
=
\mathrm{i}\mu_\xi\,\theta-\frac{\sigma_\xi^2}{2}\theta^2+O(\theta^3),
\qquad \theta\to 0,
$$
从而
$$
\log\frac{|\Lambda(\theta)|}{\lambda}
=
-\frac{\sigma_\xi^2}{2}\theta^2+O(\theta^3).
$$
特别地，当 \(\theta=2\pi/p\) 时（\(p\) 为奇素数）有
$$
1-\frac{\rho_{3,3,p}}{\lambda}
=
\frac{\sigma_\xi^2}{2}\left(\frac{2\pi}{p}\right)^2
+
O\!\left(p^{-3}\right),
$$
其中 \(\rho_{3,3,p}\) 为 \( ((3,3,p))\) 的最坏扭曲谱半径（且在可复现输出中由纯碰撞扭曲达到）。
因此二次律常数的本体为
$$
\kappa_\infty:=\lim_{p\to\infty}\frac{1-\rho_{3,3,p}/\lambda}{(2\pi/p)^2}
=\frac{\sigma_\xi^2}{2}.
$$
并且 \(\sigma_\xi^2\) 可由实参倾斜压力
\(
P(t):=\log\rho(M(1,1,e^t))
\)
给出：\(\sigma_\xi^2=P''(0)\)（与命题 \ref{prop:kernel-clt-ldp} 的有限状态 Markov 压力导数口径一致）。
\end{proposition}

\begin{remark}[选择律判别量 \texorpdfstring{$g(p)$}{g(p)}：把“模数是否特殊”变成可证伪结论]\label{rem:arity-335-selection-law-gp}
定义无量纲判别量
$$
g(p):=\left(1-\frac{\rho_{3,3,p}}{\lambda}\right)p^2.
$$
若“近共边界的一般小角微扰”机制成立，则 \(g(p)\) 应随 \(p\) 近似保持常数（仅作缓慢漂移）；若“\(p=5\) 的特异共振”机制成立，则 \(g(5)\) 将显著偏离其它素数的 \(g(p)\)。

对 \(p\in\{5,7,11,13\}\) 的谱扫描（仅计算谱半径，不依赖常数张量）由脚本
\texttt{scripts/exp\_arity\_335\_n2\_selection\_law\_primes.py}
生成并输出为可 \verb|\input| 的表格片段：
\begin{table}[H]
\centering
\scriptsize
\setlength{\tabcolsep}{6pt}
\caption{选择律扫描：第三轴取 $c=N_2(\gamma)\bmod p$ 时 $((3,3,p))$ 的最坏扭曲谱半径比与判别量 $g(p):=(1-\rho_{3,3,p}/\lambda)p^2$。这里 $\lambda=\varphi^2$。}
\label{tab:real_input_40_arity_335_n2_selection_law_primes}
\begin{tabular}{r r r r r r l}
\toprule
$p$ & $\rho_{3,3,p}$ & $\rho_{3,3,p}/\lambda$ & $1-\rho/\lambda$ & $g(p)$ & $\kappa_p$ & argmax $j$\\
\midrule
5 & 2.488758321093 & 0.950621088873 & 0.049378911127 & 1.234473 & 0.031270 & $(0,0,1),\ (0,0,4)$\\
7 & 2.549484728399 & 0.973816512450 & 0.026183487550 & 1.282991 & 0.032499 & $(0,0,1),\ (0,0,6)$\\
11 & 2.589669069543 & 0.989165564951 & 0.010834435049 & 1.310967 & 0.033207 & $(0,0,1),\ (0,0,10)$\\
13 & 2.597644541589 & 0.992211924196 & 0.007788075804 & 1.316185 & 0.033339 & $(0,0,1),\ (0,0,12)$\\
\bottomrule
\end{tabular}
\end{table}

其中同时给出达到最坏谱半径的角色索引 \(j\)（用于检验“最坏扭曲是否仍为纯碰撞扭曲 \((0,0,\pm1)\)”）。
从表中可见：在该扫描范围内 \(g(p)\) 呈现缓慢漂移并向常数量级收敛；等价地，二次律常数 \(\kappa_p:=\bigl(1-\rho_{3,3,p}/\lambda\bigr)/\bigl((2\pi/p)^2\bigr)\) 随 \(p\) 增大单调逼近其极限 \(\kappa_\infty=\sigma_\xi^2/2\)（命题 \ref{prop:arity-335-kappa-variance}）。
同时最坏角色始终为 \((0,0,1)\) 与 \((0,0,-1)\)（模 \(p\) 意义下），即慢模稳定地由纯碰撞扭曲产生，与“near-coboundary 小角微扰”机制一致。
\end{remark}

\begin{remark}[谐波能量指纹 \texorpdfstring{$\mathcal{E}_1$}{E1}：把“低秩可压缩”变成可解释的结构量]\label{rem:arity-335-harmonic-energy}
对固定的 \((a,b)\)，函数 \(c\mapsto C^{(3,3,5)}_{a,b,c}\) 在 \(\ZZ/5\) 上具有有限傅里叶展开；
由于实值性，非零谐波对只需 \(j=1,2\) 两个复模即可完整编码（命题 \ref{prop:arity-335-fourier-fingerprint}）。
因此“沿 \(c\) 轴的主要变化是否集中在第一谐波”可用无量纲指纹
$$
\mathcal{E}_1
:=\frac{\|A^{(1)}\|_F^2}{\|A^{(1)}\|_F^2+\|A^{(2)}\|_F^2}
$$
量化，其中 \(A^{(1)},A^{(2)}\) 为二级摘要输出中给出的 \(j=1,2\) 谐波复矩阵。
\(\mathcal{E}_1\) 越大，表示能量越集中在最低谐波，系统越接近“近旋转/near-coboundary”的低频结构；这与小角二次律（命题 \ref{prop:arity-335-kappa-variance}）和张量低秩可压缩现象（命题 \ref{prop:arity-335-low-rank}）在机制上同源。
\end{remark}

\begin{remark}[碰撞 cocycle 的共边界缺陷证书：best \texorpdfstring{$\ZZ/5$}{Z/5} 着色]\label{rem:arity-335-best-5-coloring}
把 40 状态图的每条边赋 cocycle $\xi=\mathbf{1}_{\{d=2\}}\in\ZZ/5$。若存在 $h:V\to\ZZ/5$ 使得 $h(v)-h(u)\equiv \xi(u\to v)\pmod 5$ 在全部边上成立，则第三轴扭曲可由共轭消去，因而不会降低谱半径（$\rho/\lambda=1$）。二级摘要输出给出一个可复现的“best found” 着色缺陷 $\varepsilon$（不一致边比例），作为“near-coboundary” 的组合学证书；其值与谱缺口 $1-\rho_{3,3,5}/\lambda$ 应在机制上相关。
\end{remark}

\begin{remark}[近共边界缺陷的组合测度与谱测度同阶]\label{rem:arity-335-eps-vs-gap}
令谱缺口
$$
\delta_{\mathrm{spec}}:=1-\frac{\rho_{3,3,5}}{\lambda}.
$$
二级摘要输出同时给出 best found 的共边界缺陷证书 \(\varepsilon\lesssim 0.066667\)（\(|E|=90\)）与谱缺口
\(
\delta_{\mathrm{spec}}\approx 0.049379
\)，并满足
$$
\delta_{\mathrm{spec}}<\varepsilon.
$$
该同阶锁定支持如下机制解释：\(\varepsilon\) 度量“少量缺陷边”对 cocycle 共边界性的破缺，而 \(\delta_{\mathrm{spec}}\) 度量同一破缺在 Perron--Frobenius 主模态上的指数效应；因此 \(((3,3,5))\) 的慢模主要来自少数破缺边诱导的相干相位阻塞，而非来自复杂的多维耦合。

进一步地，一个自然的可证伪猜想是存在仅依赖于图的粗结构常数 \(C>0\)（例如最大出度与最小正权重的比较常数），使得
$$
1-\frac{\rho(\omega_5)}{\rho(1)}\ \le\ C\,\varepsilon,
$$
其中 \(\rho(\omega_5)\) 表示沿第三轴（碰撞 cocycle）施加非平凡 \(\ZZ/5\) 角色时的最坏谱半径（即 \(\rho_{3,3,5}\)），而 \(\rho(1)=\lambda\)。
\end{remark}

\begin{remark}[从“最佳着色”得到的可证明上界：\texorpdfstring{$\kappa_\infty\le \varepsilon_\pi^{(2)}/2$}{kappa <= eps/2}]\label{rem:arity-335-eps-bound}
令 \(\xi=\mathbf{1}_{\{d=2\}}\in\{0,1\}\) 为碰撞 cocycle，并取一个着色 \(h:V\to\ZZ/5\)。
把边残差定义为
$$
r(e):=\xi(e)-\bigl(h(v)-h(u)\bigr)\quad\in\ZZ/5,
$$
并取其在 \(\{-2,-1,0,1,2\}\) 的最小整数提升（仍记为 \(r(e)\in\ZZ\)）。
则有分解 \(\xi=\delta h+r\)，其中共边界项 \(\delta h\) 在轨道和上望远镜消去；因此渐近方差密度满足
$$
\sigma_\xi^2=\sigma_r^2\le \mathbb{E}_\pi[r(e)^2]=:\varepsilon_\pi^{(2)},
$$
其中 \(\pi\) 为无扭曲系统的 PF 平衡边测度。
结合命题 \ref{prop:arity-335-kappa-variance} 的 \(\kappa_\infty=\sigma_\xi^2/2\)，得到可证明上界
$$
\kappa_\infty\ \le\ \varepsilon_\pi^{(2)}/2.
$$
二级摘要输出已同时给出 \(\varepsilon\)（均匀边比例缺陷）与 \(\varepsilon_\pi^{(2)}\)（平衡态二次缺陷），从而把“near-coboundary 证书”升级为对二次律常数的可审计上界。
\end{remark}

\begin{remark}[对照：第三维取 \texorpdfstring{$N_e\bmod 5$}{Ne mod 5} 的版本]\label{rem:arity-335-ne-mod5}
若把第三维改为输出计数 $N_e(\gamma)\bmod 5$，则得到另一个同型的 $((3,3,5))$ 常数张量（主项系数同为 $1/45$），其输出片段为
\verb|\input{sections/generated/tab_real_input_40_arity_dirichlet_mertens_335}|
（脚本默认第三维为 \texttt{Ne}）。该版本更偏“输出端”的同余分层；而本文主线采用 $N_2\bmod 5$ 以直接锁定碰撞触发机制。
\end{remark}

\paragraph{\texorpdfstring{$((5,5,5))$}{(5,5,5)}：更高模数的三维 Dirichlet--Mertens 常数张量（全表）}
为进一步提高同余分层的分辨率，可把三条观测量同时提升到模 $5$：
$$
(a,b,c)=\bigl(\chi(\gamma)\bmod 5,\ N_-(\gamma)\bmod 5,\ N_e(\gamma)\bmod 5\bigr)\in(\ZZ/5)^3.
$$
对应的 $5\times 5\times 5$ 常数张量（共 $125$ 个数）、以及最坏扭曲谱半径比 $\rho_{5,5,5}/\lambda$ 与达到最坏的角色索引，均由脚本
\texttt{scripts/exp\_sync\_kernel\_real\_input\_40\_arity\_3d.py}
直接生成并输出为可 \verb|\input| 的 \LaTeX\ 片段：
% AUTO-GENERATED by scripts/exp_sync_kernel_real_input_40_arity_3d.py
% (5,5,5) Dirichlet--Mertens constants tensor for (chi mod 5, N_- mod 5, N_2 mod 5).
\paragraph{$((5,5,5))$：Dirichlet--Mertens 常数张量（可复现输出）}
按 $a=\chi(\gamma)\bmod 5$ 分 5 张 $5\times 5$ 切片（行 $b=N_-(\gamma)\bmod 5$，列 $c=N_2(\gamma)\bmod 5$），常数项为：
$$
a=0:\quad
\begin{pmatrix}
\phantom{-}0.43171284 & -0.02910938 & -0.02155075 & -0.01848971 & -0.02647934\\
-0.02935655 & \phantom{-}0.00462483 & -0.02418489 & -0.01946126 & -0.02474958\\
-0.02890008 & \phantom{-}0.15998609 & -0.01039434 & -0.02452869 & -0.02465566\\
-0.02831606 & -0.01827209 & -0.00883335 & -0.02434939 & -0.02584702\\
-0.02849659 & -0.02629901 & \phantom{-}0.00234647 & -0.01959993 & -0.02708126
\end{pmatrix}
$$
$$
a=1:\quad
\begin{pmatrix}
\phantom{-}0.91393023 & -0.02952463 & -0.01459618 & -0.01663832 & -0.02631512\\
\phantom{-}0.23915402 & \phantom{-}0.06315752 & -0.02036496 & -0.02070844 & -0.02499819\\
-0.02886595 & \phantom{-}0.03503738 & -0.01633801 & -0.02367080 & -0.02435530\\
-0.02837130 & \phantom{-}0.02751077 & \phantom{-}0.00195706 & -0.02352503 & -0.02602958\\
-0.02844949 & -0.02725245 & -0.00775969 & -0.02074369 & -0.02702502
\end{pmatrix}
$$
$$
a=2:\quad
\begin{pmatrix}
\phantom{-}0.33542452 & -0.02979481 & -0.02029875 & -0.01755923 & -0.02628673\\
-0.02937679 & \phantom{-}0.03852094 & -0.01814409 & -0.01951852 & -0.02480101\\
\phantom{-}0.01326296 & \phantom{-}0.04091627 & -0.01501141 & -0.02439935 & -0.02455963\\
-0.02832468 & -0.01194828 & -0.00259995 & -0.02313734 & -0.02577661\\
-0.02849692 & -0.01378018 & -0.00351260 & -0.02005516 & -0.02712095
\end{pmatrix}
$$
$$
a=3:\quad
\begin{pmatrix}
\phantom{-}0.16750372 & -0.02586173 & -0.01873912 & -0.01738868 & -0.02636575\\
-0.02933676 & \phantom{-}0.02450776 & -0.02202171 & -0.02025543 & -0.02480081\\
-0.02888106 & \phantom{-}0.02466685 & -0.01484226 & -0.02373870 & -0.02443476\\
-0.02005746 & -0.01023436 & -0.00386092 & -0.02377248 & -0.02593914\\
-0.02845595 & -0.02540973 & -0.00705405 & -0.01999654 & -0.02693901
\end{pmatrix}
$$
$$
a=4:\quad
\begin{pmatrix}
\phantom{-}0.09125199 & -0.02925973 & -0.02041795 & -0.01773878 & -0.02632172\\
-0.02936272 & \phantom{-}0.01451379 & -0.02260332 & -0.02048184 & -0.02489511\\
-0.02886358 & \phantom{-}0.01338387 & -0.01721674 & -0.02421084 & -0.02449646\\
-0.02834206 & -0.01498177 & -0.00584472 & -0.02377267 & -0.02592098\\
-0.02665551 & -0.02490773 & -0.00918023 & -0.02056147 & -0.02705447
\end{pmatrix}
$$
\paragraph{指数级误差：最坏扭曲谱半径比（可复现输出）}
数值上
$$
\rho_{5,5,5}\approx 2.520679865768,\qquad \frac{\rho_{5,5,5}}{\lambda}\approx 0.962814033966.
$$
达到最坏谱半径的非平凡角色索引（$j=(j_1,j_2,j_3)$）为：
$$
(1,0,1),\ (4,0,4)
$$
\paragraph{一维边际：只看 $\chi\bmod 5$ 的 Dirichlet--Mertens 常数}
对 $b,c$ 求和得到 $C_a^{(5)}:=\sum_{b,c} C^{(5,5,5)}_{a,b,c}$：
$$
\bigl(+0.10971530,+0.84521482,-0.00637830,-0.25170809,-0.37394073\bigr).
$$


\begin{corollary}[纯碰撞慢模的普适性与混合变慢]\label{cor:arity-555-collision-universal}
在 $((5,5,5))$ 的可复现输出中，最坏扭曲角色仍为 $(0,0,1),(0,0,4)$，且
\[
\rho_{5,5,5}/\lambda\approx 0.962814> \rho_{3,3,5}/\lambda\approx 0.950621.
\]
与 $((3,3,3))$ 的最坏角色 $(1,2,2),(2,1,1)$ 及 $\rho_{3,3,3}/\lambda\approx 0.936686$ 对照可见：
提升第三维模数后，最慢模态稳定地退化为\textbf{纯碰撞扭曲}，同时 $\rho/\lambda$ 更接近 $1$，表明混合变慢、相关长度增加。
\end{corollary}

\paragraph{为何称为“真正的 arity 维度素数谱”}
上述输出同时给出三层可审计结构：
\begin{itemize}
  \item primitive 轨道作为“素数类对象”，满足 $p_n\sim \lambda^n/n$；
  \item $\chi(\gamma)$ 作为真正的 arity-charge，把二元碰撞的吐出/隐藏机制编码成可加电荷；
  \item Dirichlet--Mertens 常数张量（例如 \(((2,2,2))\)、\(((3,2,2))\)、\(((3,3,3))\)、\(((3,3,5))\)）给出同余类的常数漂移（Abel/Hadamard 有限部），而 $\rho/\lambda$ 给出同余类趋于等分布的指数误差。
\end{itemize}
